\documentclass[reqno]{amsart}
\usepackage{a4wide}
\usepackage{hyperref}

\AtBeginDocument{{\noindent\small
\emph{Electronic Journal of Differential Equations},
Vol. 2026 (2026), No. 07, pp. 1--14.\newline
ISSN: 1072-6691. URL: https://ejde.math.txstate.edu, DOI: 10.58997/ejde.2026.07}
\thanks{\copyright 2026. This work is licensed under a CC BY 4.0 license.}
\vspace{8mm}}

\begin{document}
\title[\hfilneg EJDE-2026/07\hfil 
Stability for fuzzy fractional delay differential equations]
{Finite-time stability for fuzzy fractional delay differential equations with generalized 
Caputo derivatives}


\author[K. Elhakimy, I. Sadouki, A. El Mfadel, S. Melliani \hfil EJDE-2026/07\hfilneg]
{Khadija Elhakimy, Ismail Sadouki, Ali El Mfadel, Said Melliani}

\address{Khadija Elhakimy \newline
Sultan Moulay Slimane University,
Laboratory of Applied Mathematics and Scientific Computing,
Beni Mellal, Morocco}
\email{elhakimy.khadija@usms.ac.ma}

\address{Ismail Sadouki \newline
Sultan Moulay Slimane University,
Laboratory of Applied Mathematics and Scientific Computing,
Beni Mellal, Morocco}
\email{ismail.sadouki@usms.ma}

\address{Ali El Mfadel \newline
Sultan Moulay Slimane University,
Laboratory of Applied Mathematics and Scientific Computing,
Beni Mellal, Morocco}
\email{a.elmfadel@usms.ma}

\address{Said Melliani \newline
Sultan Moulay Slimane University,
Laboratory of Applied Mathematics and Scientific Computing,
Beni Mellal, Morocco}
\email{s.melliani@usms.ma}

\thanks{Submitted September 15, 2025. Published January 19, 2026.}
\subjclass[2020]{26A33, 37C75, 34A08, 93D40}
\keywords{Fuzzy fractional differential equations; Riemann-Liouville derivative;
\hfill\break\indent generalized Caputo fractional derivative; finite-time stability}

\begin{abstract}
This work is devoted to the analysis of fuzzy fractional delay differential equations
(FDDEs) governed by the generalized Caputo fractional derivative (GCFD).
By combining stepwise approximation methods with suitable Gronwall-type inequalities,
we establish the existence and uniqueness of solutions. Furthermore, we
derive explicit criteria that guarantee the finite-time stability.
The theoretical contributions are illustrated with numerical simulations,
confirm the analytical findings and demonstrate the effectiveness
of the proposed framework in capturing the dynamical behavior of fuzzy fractional
delay models.
\end{abstract}

\maketitle
\numberwithin{equation}{section}
\newtheorem{theorem}{Theorem}[section]
\newtheorem{lemma}[theorem]{Lemma}
\newtheorem{definition}[theorem]{Definition}
\newtheorem{remark}[theorem]{Remark}
\allowdisplaybreaks

\section{Introduction}

Over the past ten years, fuzzy fractional analysis and fuzzy fractional differential
equations (FFDEs) have gained considerable attention for modeling dynamical systems with
imprecise data. Pioneering studies by Allahviranloo et al.\ \cite{allah2014}
explored explicit solutions and existence-uniqueness for FFDEs under
Riemann-Liouville and Caputo fuzzy derivatives.
Lupulescu \cite{lupu} formulated a fractional calculus framework for interval-valued
functions, later applied to FFDEs. Fard and Salehi \cite{sal} studied fuzzy variational
problems using the Caputo derivative, while Hoa \cite{hoa2020a} extended Caputo fuzzy
derivative theory to orders in $(1,2)$, establishing solution uniqueness.
Hoa and Vu \cite{hoa2019b} investigated qualitative properties of implicit FFDEs
with Caputo, Riemann-Liouville, and Hadamard derivatives. Long et al.\
\cite{long} introduced Caputo fuzzy partial derivatives for multivariable functions,
addressing fuzzy fractional partial differential equations. Research on FFDEs has
largely centered on operators such as the Riemann-Liouville, Caputo, Caputo-Hadamard,
and Caputo-Katugampola derivatives within the generalized Hukuhara derivative framework.
The choice of fractional derivative is guided by the alignment of experimental data
with theoretical models, necessitating more flexible fractional operators that
interpolate between classical forms. Samko \cite{samk} introduced the generalized
fractional derivative, defined for a real-valued function $ w$ with respect
to a function $\psi$. Almeida \cite{alm} advanced this by proposing the GCFD ,
which unifies various fractional derivative types. Recent studies have leveraged the
GCFD Baleanu et al.\ \cite{bal} investigated the well-posedness of terminal value
problems for nonlinear systems, Wu et al.\ \cite{shiri} applied it to linear
systems with non-homogeneous terms, and Luo et al.\ \cite{luo} used it to study
the qualitative behavior of solutions to fractional uncertain differential equations
driven by the Liu process. Furthermore, Fan et al.\ \cite{Fan} laid the groundwork for
general fractional calculus in the space of interval-valued functions, providing a
theoretical basis for future advancements in FFDEs.
Continued exploration of the references cited in these works offers valuable
insights into the evolving landscape of generalized fractional operators and their
applications. The reader may refer to the recent works
\cite{abbas,ahm1,akd,ali1,ali2,ali3,xi} for more details.

Stability analysis is a cornerstone of control theory, especially in modeling dynamic
systems. FTS, which ensures that system trajectories remain bounded within a specified
time interval, has garnered significant interest due to its practical importance.
Unlike asymptotic stability, which addresses long-term behavior, FTS focuses on state
boundedness over a finite duration, making it ideal for time-critical applications.
A system is deemed finite-time stable if its states stay within prescribed bounds
during a given time window. Time delays, which account for memory effects in
real-world systems, are often critical to include in dynamic models.
In \cite{laz2006}, Lazarević pioneered the study of FTS in fractional-order delay
differential systems by establishing sufficient conditions. Since then,
the field has advanced rapidly in both theoretical and applied domains, utilizing
various fractional derivative definitions. The literature on FTS in fractional delay
systems splits into two primary approaches Chen et al. in \cite{che} use H\"older-type
inequalities and analytical methods to derive straightforward sufficient conditions
for FTS in nonlinear fractional systems with delays, while Zhang et al.\ \cite{zhan}
focus on systems with variable coefficients. Additionally, Naifar et al.\
\cite{naif} employed the fractional Gronwall inequality to study FTS in linear
time-delay systems, an approach extended to nonlinear systems with time-varying delays
by Phat et al.\ \cite{phat}. Extensive investigations into FTS across fractional
differential systems have been conducted, yet the case of FDDEs
 with time-varying delays, approached through the fuzzy framework and the GCFD,
remains largely uninvestigated. This important literature gap forms the central
motivation for the current research.

In this article, we aim to fill this gap by investigating the finite-time stability
(FTS) of fuzzy fractional delay differential equations (FDDEs) within the generalized
Caputo fractional derivative (GCFD) framework.
Specifically, we consider the system
\begin{equation} \label{sys11}
\begin{gathered}
{}^C \mathcal{D}_{0^{+}}^{\beta ; \psi} w(\nu )=\mathcal{H}(\nu , w(\nu ),
 w(\nu - r)),\quad \nu \in \mathcal{J}=[0, T], \\
 w(\nu )=\varrho(\nu ) \in \mathfrak{F}, \quad \nu \in[- r, 0],
\end{gathered}
\end{equation}
where
\begin{itemize}
\item $\beta \in (0,1)$ and $r> 0$.

\item The interval $\mathcal{J}$ is partitioned as $\mathcal{J}
= \cup_{k=0}^{n} [k r, (k+1) r] \cup [(n+1) r, T]$.

\item $\mathfrak{F}$ the space of fuzzy numbers.

\item The function
$\mathcal{H} : \mathcal{J} \times \mathfrak{F} \times \mathfrak{F} \to \mathfrak{F}$,
which is jointly continuous and satisfies the condition $\mathcal{H}(\nu , \hat{0},
\hat{0}) = \hat{0}$.

\item The operator ${}^C \mathcal{D}_{0^{+}}^{\beta;\psi} w$ represents the
 GCFD of the fuzzy-valued function $ w$ with respect to the lower limit $0^+$.

\item The initial function $\varrho(\nu )$ is assumed to be continuous on the
interval $[-r, 0]$.

\item The function $\psi : \mathcal{J} \to \mathbb{R}$ belongs to the class
$\mathcal{K}$, which consists of functions $\psi \in C^1$ that are strictly
increasing, positive, and satisfy $\psi'(\nu ) \neq 0$ for all $\nu \in (0, T)$.
\end{itemize}

Our work provides new insights into the interplay between fuzziness, memory effects,
and delay phenomena in fractional dynamical systems.
The structure of the paper is as follows:
Section 2 covers the essential groundwork, including notations and fundamental concepts
needed for the main results, and presents the explicit solution for \eqref{sys11}.
Section 3 investigates the existence and uniqueness of solutions, proposing two
separate methods to establish the FTS of \eqref{sys11} with appropriate sufficient
conditions.
Section 4 offers a range of examples that substantiate the theoretical findings
and illustrate the practical relevance of the suggested approaches.

\section{Preliminaries}

 The following concepts are fundamental to understanding the structure and
 operations within $\mathfrak{F}$ we direct readers to \cite{bed}.
 
For any fuzzy number \( z \in \mathfrak{F} \), the \(\zeta\)-cuts set for
\(\zeta \in (0,1]\) is defined as
 \[
 [z]^\zeta = \{ x_0 \in \mathbb{R} \mid z(x_0) \geq \zeta \},
 \]
and the 0-cuts set is given by
\[
 [z]^0 = \overline{\{ x_0 \in \mathbb{R} \mid z(x_0) > 0 \}}.
\]
Each \([z]^\zeta\) forms a closed interval \([\underline{z}(\zeta), \overline{z}(\zeta)]\). The zero fuzzy number \(\hat{0} \in \mathfrak{F}\) is defined by
\[
 \hat{0}(x_0) = \begin{cases}
 1 & \text{if } x_0 = 0, \\
 0 & \text{otherwise}.
 \end{cases}
\]

According to Zadeh's extension principle, addition and scalar multiplication 
in $\mathfrak{F}$ are defined for all $\zeta \in [0,1]$ as
\[
 [z_1 + z_2]^\zeta = [z_1]^\zeta + [z_2]^\zeta, \quad 
 [\lambda z_1]^\zeta = \lambda [z_1]^\zeta,
\]
 where
\[
 [z_1]^\zeta + [z_2]^\zeta = [\underline{z_1}(\zeta) + \underline{z_2}(\zeta), \overline{z_1}(\zeta) + \overline{z_2}(\zeta)],
 \]
and
\begin{equation} \label{eq21}
 \lambda [z_1]^\zeta = \begin{cases}
 [\lambda \underline{z_1}(\zeta), \lambda \overline{z_1}(\zeta)]
 & \text{if }   \lambda \geq 0, \\
 [\lambda \overline{z_1}(\zeta), \lambda \underline{z_1}(\zeta)]
 & \text{if } \lambda < 0.
 \end{cases}
\end{equation}
 Alternatively, these operations can be expressed as
\[
 [z_1]^\zeta + [z_2]^\zeta = \{x_0 + y_0 : x_0 \in [z_1]^\zeta, y_0 \in [z_2]^\zeta\},
 \quad \lambda [z_1]^\zeta = \{\lambda x_0 : x_0 \in [z_1]^\zeta\}.
\]

 The diameter of the \(\zeta\)-cuts set of \(z\) is
 \[
 d([z]^\zeta) = \overline{z}(\zeta) - \underline{z}(\zeta), \quad \zeta \in [0,1].
 \]

The Hukuhara difference \(z_1 \ominus z_2\) exists if there exists
 \(z_3 \in \mathfrak{F}\) such that \(z_1 = z_2 + z_3\).
The generalized Hukuhara difference (gH-difference) extends this 
concept and is defined as
 \[\label{eq22}
 z_1 \ominus_{gH} z_2 = z_3 \Leftrightarrow
 \begin{cases}
 \text{(i) } z_1 = z_2 + z_3, \quad \text{or} \\
 \text{(ii) } z_2 = z_1 + (-1)z_3.
 \end{cases}\tag{2.2}
 \]
Case (i) applies when \(d([z_1]^\zeta) \geq d([z_2]^\zeta)\), and case (ii) when 
\(d([z_1]^\zeta) \leq d([z_2]^\zeta)\).

 The metric on \(\mathfrak{F}\) is defined via the Hausdorff distance
\[
 D_0[z_1, z_2] = \sup_{\zeta \in [0,1]} H([z_1]^\zeta, [z_2]^\zeta),
\]
where
$$
 H([z_1]^\zeta, [z_{2}]^\zeta) 
 = \max \{ | \underline{z_1}(\zeta) - \underline{z_2}(\zeta) |,
 | \overline{z_1}(\zeta) - \overline{z_2}(\zeta) | \}.
$$

 A function \( w : [a,b] = \mathfrak{I} \to \mathfrak{F}\) with 
fuzzy values is defined as \(d\)-monotone if, for any fixed \(\zeta \in [0,1]\), 
the function \(\nu \mapsto d([ w(\nu )]^\zeta)\) is monotonic across \(\mathfrak{I}\).

\begin{definition}[\cite{bed}] \label{def1} \rm
Consider a fuzzy-valued function \(\mathfrak{g} : (a, b) \to \mathfrak{F}\). 
The gH-derivative of \(\mathfrak{g}\) at a point \(\nu \in (a, b)\) is expressed as
\begin{equation} \label{e2.3}
\mathfrak{g}'(\nu ) = \lim_{h \to 0} \frac{\mathfrak{g}(\nu + h)
\ominus_{gH} \mathfrak{g}(\nu )}{h}.
\end{equation}
The function \(\mathfrak{g}\) is deemed gH-differentiable at \(\nu \) if the limit
exists and belongs to \(\mathfrak{F}\).
\end{definition}

Let \( C^1(\mathfrak{I}, \mathfrak{F}) \) denote the space of all fuzzy functions 
\( w : \mathfrak{I} \to \mathfrak{F} \) that are generalized Hukuhara differentiable 
(gH-differentiable) on \( \mathfrak{I} \) and whose gH-derivative \( w' \) is 
continuous, i.e., \( w' \in C(\mathfrak{I}, \mathfrak{F}) \).

 For \(\mathfrak{g},\mathfrak{f} \in C(\mathfrak{I}, \mathfrak{F})\), we define
 \[
 D_*[\mathfrak{g}, \mathfrak{f}] 
 = \sup_{\nu \in \mathfrak{I}} D_0[\mathfrak{g}(\nu ),\mathfrak{f}(\nu )].
 \]
 Under \((C(\mathfrak{I}, \mathfrak{F}),D_*)\) is a complete metric space.

 A function \(\mathcal{H}: \mathfrak{I} \times C(\mathfrak{I}, \mathfrak{F}) 
 \to \mathfrak{F}\) is jointly continuous at \((\nu _0, w_0)\) if
 \[
 \forall \epsilon > 0, \exists \delta > 0 \text{ such that } | \nu -\nu _0| + D_0[ w, w_0] < \delta \implies D_0[\mathcal{H}(\nu , w), \mathcal{H}(\nu _0, w_0)] < \epsilon.
 \]
The notation \( C_{\mathbb{J}}(\mathfrak{I}, \mathfrak{F}) \) refers to the 
space of jointly continuous functions.

\begin{definition}[\cite{alm}] \label{def2} \rm
Consider \(\psi \in \mathcal{K}\), where \(\mathcal{K}\) consists of strictly 
increasing, positive, \(C^1\)-functions on \(\mathfrak{I}\) with non-vanishing 
derivatives on \((a, b)\). The generalized Riemann-Liouville fuzzy fractional integral 
of order \(\beta \in (0,1)\) for a fuzzy-valued function 
\(\mathfrak{g} : \mathfrak{I} \to \mathfrak{F}\) is given by
\[
{}^{c}\mathcal{I}_{a^{+}}^{\beta ; \psi} \mathfrak{g}(\nu ) 
= \frac{1}{\Gamma(\beta)} \int_{a}^{\nu } \psi'(s) \mathfrak{D}_{\nu ,s}^{\beta - 1}
 \mathfrak{g}(s) \, ds, \quad \text{for } \nu > a,
\]
where $\mathfrak{D}_{\nu ,s}=\psi(\nu )-\psi(s)$, and \(\Gamma(\cdot)\) denotes the 
Gamma function.
\end{definition}

Based on the GCFD framework introduced for real-valued functions in \cite{alm,bal} 
and extended to interval-valued functions in \cite{Fan}, we propose an extension of the 
generalized fractional derivative definition to the context of fuzzy-valued functions.

\begin{definition}\label{def3} \rm
Suppose \(\psi \in \mathcal{K}\), \(\beta \in (0,1)\), and let 
\(\mathfrak{g} : \mathfrak{I} \to \mathfrak{F}\) be a fuzzy value function where 
\({}^{c}\mathcal{I}_{a^{+}}^{1-\beta ; \psi} \mathfrak{g}(\nu )\) is gH-differentiable 
in \(\mathfrak{I}\). The generalized Riemann--Liouville fractional derivative of order 
\(\beta\) is given by
\[
{}^{RL} \mathcal{D}_{a^{+}}^{\beta ; \psi} \mathfrak{g}(\nu ) 
= \frac{1}{\psi'(\nu )} ( {}^{c} \mathcal{I}_{a^{+}}^{1-\beta ; \psi} \mathfrak{g}(\nu ) )',
\]
equivalently written as
\[
{}^{RL} \mathcal{D}_{a^{+}}^{\beta ; \psi} \mathfrak{g}(\nu )
 = \frac{1}{\psi'(\nu )} \cdot \frac{1}{\Gamma(1 - \beta)} ( \int_{a}^{\nu } \psi'(s) \mathfrak{D}_{\nu ,s}^{-\beta} \mathfrak{g}(s) \, ds )', \quad \text{for } \nu > a. \tag{2.5}
\]
\end{definition}

Based on Definition \ref{def3}, the GCFD of the fuzzy-valued function $\mathfrak{g}$ 
is defined in the weak sense as follows.


\begin{definition}\label{def4} \rm
Let \(\psi \in \mathcal{K}\), \(\beta \in (0,1)\), and let 
\(\mathfrak{g} : \mathfrak{I} \to \mathfrak{F}\) be a fuzzy-valued function. If the 
generalized Riemann-Liouville fractional derivative 
\({}^{RL} \mathcal{D}_{a^{+}}^{\beta ; \psi} \mathfrak{g}(\nu )\) exists on 
\(\mathfrak{I}\), then the GCFD of \(\mathfrak{g}\) is defined in the weak sense by
$$
{}^C \mathcal{D}_{a^{+}}^{\beta ; \psi} \mathfrak{g}(\nu ) 
= {}^{RL} \mathcal{D}_{a^{+}}^{\beta ; \psi} [ \mathfrak{g}(\cdot) \ominus_{gH} 
\mathfrak{g}(a) ],
$$
which can be equivalently expressed as
\[
{}^C \mathcal{D}_{a^{+}}^{\beta ; \psi} \mathfrak{g}(\nu ) 
= \frac{1}{\psi'(\nu )} \cdot \frac{1}{\Gamma(1 - \beta)} 
\Big( \int_{a}^{\nu } \psi'(s) \mathfrak{D}_{\nu ,s}^{-\beta} [ \mathfrak{g}(s)
 \ominus_{gH} \mathfrak{g}(a) ] ds \Big)', \quad \nu > a. \tag{2.6}
\]
\end{definition}

\begin{theorem}[\cite{hoa2018}] \label{thm2.1}
Let \(\psi \in \mathcal{K}\) and \(\beta \in (0,1)\). 
If \(\mathfrak{g} \in C^1(\mathfrak{I}, \mathfrak{F})\) is a d-monotone fuzzy function, 
then the GCFD of \(\mathfrak{g}\) is given by
\begin{equation}
\label{e2.7}
{}^C \mathcal{D}_{a^{+}}^{\beta ; \psi} \mathfrak{g}(\nu )
= {}^{c} \mathcal{I}_{a^{+}}^{1-\beta ; \psi} ( \frac{1}{\psi'(\nu )}
\mathfrak{g}'(\nu ) ) = \frac{1}{\Gamma(1 - \beta)} \int_{a}^{\nu }
\mathfrak{D}_{\nu ,s}^{-\beta} \mathfrak{g}'(s) \, ds, \quad \nu > a,
\end{equation}
where the term \(\frac{1}{\psi'(\nu )} \mathfrak{g}'(\nu )\) denotes the scalar
multiplication of a real-valued function with a fuzzy-valued function.
\end{theorem}

\begin{remark} \label{rmk2.2} \rm
Equation \eqref{e2.7} defines \(^C \mathcal{D}_{a^{+}}^{\beta ; \psi} \mathfrak{g}(\nu )\)
as the strong-sense generalized Caputo fractional derivative. 
Moreover, strong GCFD differentiability entails weak GCFD differentiability for all 
\(\mathfrak{g} \in C^1(\mathfrak{I}, \mathfrak{F})\).
\end{remark}

Using a proof of technique analogous to that employed in 
\cite[Prop. 1 and Corollary 1]{hoa2018}, we can establish the following results.

\begin{theorem}\label{thm2}
Let \(\psi \in \mathcal{K}\), \(\beta \in (0,1)\), and let 
\(\mathfrak{g} \in C(\mathfrak{I}, \mathfrak{F})\) be a d-monotone fuzzy function on 
the interval \(\mathfrak{I}\). Then, the following properties hold
\begin{enumerate}
\item ${}^{c} \mathcal{I}_{a^{+}}^{\beta ; \psi} ( {}^C \mathcal{D}_{a^{+}}^{\beta ; 
\psi} \mathfrak{g}(\nu ) ) = \mathfrak{g}(\nu ) \ominus_{gH} \mathfrak{g}(a)$;

\item ${}^{RL} \mathcal{D}_{a^{+}}^{\beta ; \psi} ( {}^{c} 
\mathcal{I}_{a^{+}}^{\beta ; \psi} \mathfrak{g}(\nu ) ) = \mathfrak{g}(\nu )$;

\item ${}^C \mathcal{D}_{a^{+}}^{\beta ; \psi} ( {}^{c} \mathcal{I}_{a^{+}}^{\beta ; 
\psi} \mathfrak{g}(\nu ) ) = \mathfrak{g}(\nu )$.
\end{enumerate}
\end{theorem}

We set $\mathbb{G}(\nu )=\mathcal{H}(\nu , w(\nu ), w(\nu - r ))$. Similar to the 
\cite[Theorem 3]{hoa2018}, we also obtain the lemma below.

\begin{lemma}\label{lem1}
Assume that \(\mathbb{G} : \mathcal{J} \to \mathfrak{F}\) is jointly continuous. 
Then, a \(d\)-monotone function \( w \in C([- r, T], \mathfrak{F})\) solves problem 
\eqref{sys11} if and only if it satisfies
\begin{equation} \label{sys28}
\begin{gathered}
 w(\nu )=\varrho(\nu ), \quad \nu \in[- r, 0]  \\
 w(\nu ) \ominus_{g H} \varrho(0)=\frac{1}{\Gamma(\beta)} \int_0^{\nu }
 \psi'(s)\mathfrak{D}_{\nu ,s}^{\beta-1} \mathbb{G}(s)\,ds, \quad \nu \in\mathcal{J}
\end{gathered}
\end{equation}
assuming the right-hand side integral exhibits \(d\)-monotonicity on \((0, T]\).
\end{lemma}

\begin{proof}
Consider \( w \in C([- r, T], \mathfrak{F})\) as a \(d\)-monotone solution to 
problem \eqref{sys11} on \(\mathcal{J}\). For \(\nu \in [- r, 0]\), it holds that 
\( w(\nu ) = \varrho(\nu )\). Then, for \(\nu \in \mathcal{J}\), applying the operator 
\({}^c \mathcal{I}_{0^{+}}^{\beta ; \psi}\) to both sides of \eqref{sys11} and 
according to property (i) of Theorem \ref{thm2}, we derive
\begin{equation}\label{eq29}
 w(\nu ) \ominus_{g H} \varrho(0)=\frac{1}{\Gamma(\beta)}
  \int_0^{\nu } \psi'(s)\mathfrak{D}_{\nu ,s}^{\beta-1} \mathbb{G}(s)\,ds,
  \quad \nu \in\mathcal{J} 
\end{equation}
Consequently, \( w\) is a solution to \eqref{sys28}. It is also known that the \(d\)-
increasing property of \eqref{eq29}'s left-hand side over \(\mathcal{J}\) ensures the 
same \(d\)-increasing behavior for its right-hand side over \(\mathcal{J}\).

Conversely, suppose that \( w \in C([- r, T], \mathfrak{F})\) is a \(d\)-monotone 
fuzzy function that satisfies equation \eqref{sys28}, and that the integral on the 
right-hand side is \(d\)-increasing on \(\mathcal{J}\). It is straightforward to verify 
that \( w(\nu ) = \varrho(\nu )\) for all \(\nu \in [- r, 0]\). 
Since the fuzzy function \(\mathbb{G}\) is jointly continuous, the fractional integral 
\({}^c \mathcal{I}_{0^{+}}^{\beta ; \psi} \mathbb{G}(\nu )\) is also continuous, 
and satisfies \({}^c \mathcal{I}_{0^{+}}^{\beta ; \psi} \mathbb{G}(0) = \hat{0}\). 
Consequently, we have \( w(0) = \varrho(0)\).

In addition, since the mapping 
\(\nu \mapsto {}^c \mathcal{I}_{0^{+}}^{\beta ; \psi} \mathbb{G}(\nu )\) is 
\(d\)-increasing on \(\mathcal{J}\), applying the generalized Riemann-Liouville 
derivative \({}^{RL} \mathcal{D}_{0^{+}}^{\beta ; \psi}\) to equation \eqref{sys28} 
and invoking part (ii) of Theorem \ref{thm2} along with Definition \ref{def4}, 
we obtain
\[
{}^C \mathcal{D}_{0^{+}}^{\beta ; \psi} w(\nu ) = \mathbb{G}(\nu ),
\]
which confirms that equation \eqref{sys11} is satisfied.
\end{proof}

\begin{definition} \label{def5} \rm
Let \( w : [- r, T] \to \mathfrak{F}\) be a function that solves problem \eqref{sys11}. 
It is a solution if \( w \in C([- r, T], \mathfrak{F})\), adheres to 
\( w(\nu ) = \varrho(\nu )\) for \(\nu \in [- r, 0]\), and satisfies 
\({}^C \mathcal{D}_{0^{+}}^{\beta ; \psi} w(\nu ) = \mathbb{G}(\nu )\) on 
\(\mathcal{J}\). We classify \( w\) as \(d\)-monotone if it is either \(d\)-increasing
or \(d\)-decreasing throughout \(\mathcal{J}\).
\end{definition}

\begin{remark}\label{rq22} \rm
According to equation \eqref{eq22}, which defines the generalized Hukuhara difference, 
the integral equation \eqref{sys28} can be recast as follows

If \(d([ w(\nu )]^{\zeta}) \geq d([\varrho(0)]^{\zeta})\) for all 
\(\nu \in \mathcal{J}\) and every \(\zeta \in [0,1]\), so that 
\(\nu \mapsto d([ w(\nu )]^{\zeta})\) is nondecreasing on \(\mathcal{J}\), 
then \eqref{sys28} is expressed as
\begin{equation} \label{sys210}
\begin{gathered}
 w(\nu ) = \varrho(\nu ), \quad \nu \in [- r, 0] \tag{2.10} \\
 w(\nu ) = \varrho(0) + \frac{1}{\Gamma(\beta)} \int_0^{\nu } \psi'(s)
  \mathfrak{D}_{\nu ,s}^{\beta-1} \mathbb{G}(s) \, ds, \quad \nu \in \mathcal{J}
\end{gathered}
\end{equation}

If \(d([ w(\nu )]^{\zeta}) \leq d([\varrho(0)]^{\zeta})\) for all 
\(\nu \in \mathcal{J}\) and every \(\zeta \in [0,1]\), indicating that 
\(\nu \mapsto d([ w(\nu )]^{\zeta})\) is nonincreasing on \(\mathcal{J}\), 
then \eqref{sys28} becomes
\begin{equation} \label{sys211}
\begin{gathered}
 w(\nu ) = \varrho(\nu ), \quad \nu \in [- r, 0] \tag{2.11} \\
 w(\nu ) = \varrho(0) \ominus \frac{(-1)}{\Gamma(\beta)} \int_0^{\nu }
  \psi'(s) \mathfrak{D}_{\nu ,s}^{\beta-1} \mathbb{G}(s) \, ds, \quad
  \nu \in \mathcal{J}
\end{gathered}
\end{equation}

Assuming the Hukuhara difference in \eqref{sys211} exists on \(\mathcal{J}\), 
the two cases of the generalized Hukuhara difference imply that \eqref{sys28} acts 
as a general form encompassing one of these integral equations.

Let \(\mathcal{X}^{(i)}\) be the collection of fuzzy functions 
\( w \in C([- r, T], \mathfrak{F})\) such that 
\({}^C \mathcal{D}_{0^{+}}^{\beta ; \psi} w \in C(\mathcal{J}, \mathfrak{F})\).
Further, let \(\mathcal{X}^{(ii)}\) denote the subset of \(\mathcal{X}^{(i)}\) 
where the Hukuhara difference
\[
 w(\nu ) = \varrho(0) \ominus \frac{(-1)}{\Gamma(\beta)} \int_0^{\nu } \psi'(s)
 \mathfrak{D}_{\nu ,s}^{\beta-1} \mathbb{G}(s) \, ds,
\]
is well-defined for \(\nu \in \mathcal{J}\) and \(\psi \in \mathcal{K}\).
\end{remark}

\begin{theorem}[Schaefer's fixed point theorem, \cite{smar}] \label{theo3}
Let the operator $\mathcal{Q}: \mathcal{X} \to \mathcal{X}$ be completely continuous, 
where $\mathcal{X}$ is a Banach space. We denote
$$
\mathfrak{P}(\mathcal{Q})=\{v \in \mathcal{X}
\mid v=\lambda \mathcal{Q}(v) \text { for some } 0<\lambda<1\} .
$$
Then, if $\mathfrak{P}(\mathcal{Q})$ is bounded, $\mathcal{Q}$ has at least one fixed 
point.
\end{theorem}


To analyze the existence, uniqueness, and FTS of solutions to problem \eqref{sys11}, 
we use the following assumptions:
\begin{itemize}
\item[(A1)] The fuzzy mapping \(\mathcal{H} : \mathcal{J} \times \mathfrak{F} 
\times \mathfrak{F} \to \mathfrak{F}\) is jointly continuous.

\item[(A2)] A constant \(L > 0\) exists such that
$$
D_0[\mathcal{H}(\nu , z_{1}, z_{2}), \mathcal{H}(\nu , w_{1}, w_{2})] \leq L(D_0[z_{1}(\nu ), w_{1}(\nu )]+D_0[z_{2}(\nu - r), w_{2}(\nu - r)]).
$$

\item[(A3)] There exist positive constants $\hat{M}_{1}$ and $\hat{M}_{2} \in(0,1)$ 
such that
$$
D_0[\mathcal{H}(\nu , z_{1}, z_{2}), \hat{0}] \leq \hat{M}_{1} D_0[z_{1}(\nu ), \hat{0}]+\hat{M}_{2} D_0[z_{2}(\nu - r), \hat{0}], \forall \nu \in\mathcal{J}.
$$
\end{itemize}

\section{Existence and uniqueness results}
We establish the existence of solutions to problem \eqref{sys11}. 
To achieve this, we define two operators as follows.

To secure a \(d\)-increasing solution, we define the operator 
\(\mathcal{Q} : \mathcal{X}^{(i)} \to \mathcal{X}^{(i)}\) as
\begin{equation} \label{sys31}
(\mathcal{Q} w)(\nu )=\begin{cases}
\varrho(\nu ), & \nu \in[- r, 0] \tag{3.1}\\
\varrho(0)+\frac{1}{\Gamma(\beta)} \int_0^{\nu }
\psi'(s)\mathfrak{D}_{\nu ,s}^{\beta-1} \mathbb{G}(s)\,ds, & \nu \in\mathcal{J}
\end{cases}
\end{equation}

For the existence of the $d$-decreasing solution, we consider the operator
 $\mathcal{Q}: \mathcal{X}^{(ii)} \to \mathcal{X}^{(ii)}$ given by
\begin{equation} \label{sys32}
(\mathcal{Q} w)(\nu )=\begin{cases}
\varrho(\nu ), & \nu \in[- r, 0]  \\
\varrho(0) \ominus \frac{(-1)}{\Gamma(\beta)} \int_0^{\nu }
\psi'(s)\mathfrak{D}_{\nu ,s}^{\beta-1} \mathbb{G}(s)\,ds, & \nu \in\mathcal{J}
\end{cases}
\end{equation}
where $\mathbb{G}(\nu )=\mathcal{H}(\nu , w(\nu ), w(\nu - r))$.

Let $T^{*} \in (0,T]$ such that
$\lambda(\hat{M}_{1}+\hat{M}_{2})\mathfrak{D}_{T^{*},0}^{\beta}<\Gamma(\beta+1)$, 
where $\beta \in(0,1)$ and $\lambda \in(0,1)$.

\begin{theorem} \label{thm3.1}
If {\rm (A1)--(A3)} are satisfied, then a solution to problem \eqref{sys11} 
exists on \([0, T^*]\) for every case under study.
\end{theorem}

\begin{proof}
Since the proof strategy is the same for both cases in Remark \ref{rq22}, 
we focus on a representative one. 
The existence of a solution to \eqref{sys11} in (A1)--(A3) follows from 
Schaefer's fixed point theorem.

We reexamine the operator $\mathcal{Q}$ from \eqref{sys32} and verify that it is 
correctly defined by confirming 
$\mathcal{Q}(\mathcal{X}^{(ii)}) \subseteq \mathcal{X}^{(ii)}$. 
This conclusion follows immediately from the continuity of $\mathbb{G}$, which ensures 
the continuity of $(\mathcal{Q} w)(\nu )$ for $\nu \in [- r, T]$.
Applying \({}^C \mathcal{D}_{0^{+}}^{\beta ; \psi}\) to both sides of \eqref{sys32} gives
\[
{}^C \mathcal{D}_{0^{+}}^{\beta ; \psi}(\mathcal{Q} w)(\nu ) 
= \mathbb{G}(\nu ), \quad \forall \nu \in \mathcal{J},
\]
which implies that \({}^C \mathcal{D}_{0^{+}}^{\beta ; \psi}(\mathcal{Q} w)\) 
is continuous on \(\mathcal{J}\).
The verification of the hypotheses of Theorem \ref{theo3} will now be carried out 
in four  steps.
\smallskip

\noindent\textbf{Step 1.} The operator \(\mathcal{Q} : \mathcal{X}^{(ii)} 
\to \mathcal{X}^{(ii)}\) is continuous. To demonstrate this, take a sequence 
\(\{ w_n\}_{n \in \mathbb{N}} \subset \mathcal{X}^{(ii)}\) such that
 \( w_n \to w \in \mathcal{X}^{(ii)}\) as 
\(n \to \infty\). For \( \nu \in [- r, 0] \), it holds that
\[
\mathcal{Q}( w_n)(\nu ) = \varrho(\nu ) \to \mathcal{Q}( w)(\nu ) = \varrho(\nu ),
\]
since \(\mathcal{Q}\) is defined to yield \(\varrho(\nu )\) for all functions 
in \(\mathcal{X}^{(ii)}\) on \([- r, 0]\).
$$
D_0[(\mathcal{Q} w_{n})(\nu ),(\mathcal{Q} w)(\nu )]=0,
$$
and for any $\nu \in\mathcal{J}$, one also obtains
$$
D_0[(\mathcal{Q} w_{n})(\nu ),(\mathcal{Q} w)(\nu )]
\leq  \frac{2L}{\Gamma(\beta+1)}\mathfrak{D}_{\nu ,0}^{\beta} \sup _{\nu \in[- r, T]} D_0[ w_{n}(\nu ), w(\nu )],
$$
where we use assumption (A2). Hence, it yields
$$
D_{*}[\mathcal{Q} w_{n}, \mathcal{Q} w] 
\leq \frac{2L\mathfrak{D}_{\nu ,0}^{\beta} }{\Gamma(\beta+1)}D_{*}[ w_{n}, w], 
\quad \forall \nu \in[- r, T] .
$$
Because $ w_{n} \to w$ as $n \to \infty$, one can imply that 
$D_{*}[\mathcal{Q} w_{n}, \mathcal{Q} w] \to 0$ on $\mathcal{X}^{(ii)}$ as 
$n \to \infty$. Therefore, the operator $\mathcal{Q}$ is continuous on 
$\mathcal{X}^{(ii)}$.
\smallskip

\noindent\textbf{Step 2.} 
Consider the ball \( \mathbb{B}(\hat{0}, \rho) \subset \mathcal{X}^{(ii)} \), 
defined by \( D_0[\mathfrak{g}(\nu ), \hat{0}] \leq \rho \). 
We aim to prove that \( \mathcal{Q} \) is a bounded operator on bounded subsets, 
i.e., there exists \( \hat{\rho} > 0 \) such that 
\( D_{*}[\mathcal{Q}v, \hat{0}] \leq \hat{\rho} \). This follows from the behavior of 
\( \mathcal{Q} \) on \( [-r, 0] \) and from assumptions (A1) and (A2) on 
\( \mathcal{J} \).
\begin{align*}
D_0[(\mathcal{Q} w)(\nu ), \hat{0}]
&\leq  D_0[\varrho(0), \hat{0}]
 +\frac{L}{\Gamma(\beta)}\Big(\sup _{\nu \in[- r, T]} D_0[ w(\nu ), \hat{0}]
 +\sup _{\nu \in[- r, T]} D_0[ w(\nu - r), \hat{0}]\Big) \\
&\quad\times \int_0^{\nu } \psi'(s)\mathfrak{D}_{\nu ,s}^{\beta-1}\,ds
 +\frac{1}{\Gamma(\beta)} \int_0^{\nu } \psi'(s)\mathfrak{D}_{\nu ,s}^{\beta-1}
  D_0[\mathcal{H}(s, \hat{0}, \hat{0}), \hat{0}]\,ds \\
&\leq  \rho+ \frac{2L\mathfrak{D}_{\nu ,0}^{\beta}}{\Gamma(\beta+1)} D_{*}[ w, \hat{0}].
\end{align*}
Therefore, by choosing
$$
\hat{\rho} = \rho + \frac{2L\mathfrak{D}_{\nu ,0}^{\beta} D_{*}
[ w, \hat{0}]}{\Gamma(\beta+1)},
$$
we ensure that the operator \( \mathcal{Q} \) maps any function 
\( w \in \mathbb{B}(\hat{0}, \rho) \) into a function whose distance from 
\( \hat{0} \) remains bounded by \( \hat{\rho} \).
$$
D_{*}[\mathcal{Q} w, \hat{0}] \leq \hat{\rho}<\infty, \quad \forall \nu \in\mathcal{J} .
$$
This shows that the operator $\mathcal{Q}$ maps bounded sets into bounded sets of 
$\mathcal{X}^{(ii)}$.
\smallskip

\noindent\textbf{Step 3.} 
We now show that the operator $\mathcal{Q}$ maps bounded sets into equi-continuous 
subsets of $\mathcal{X}^{(ii)}$. To this end, consider two time points 
$\nu _1, \nu _2 \in \mathcal{J}$ such that $\nu _1 < \nu _2$, and let 
$ w \in \mathbb{B}(\hat{0}, \rho)$. Then, based on assumptions (A1) and 
(A2), one can derive that
\begin{align*}
&D_0[(\mathcal{Q} w)(\nu _{2}),(\mathcal{Q} w)(\nu _{1})] \\
&\leq \frac{1}{\Gamma(\beta)} \int_0^{\nu _{1}} \psi'(s)
 (\mathfrak{D}_{\nu _{2},s}^{\beta-1}-\mathfrak{D}_{\nu _{1},s}^{\beta-1}) 
 D_0[\mathbb{G}(s), \hat{0}]\,ds
 +\frac{1}{\Gamma(\beta)} \int_{\nu _{1}}^{\nu _{2}} \psi'(s)
 \mathfrak{D}_{\nu _{2},s}^{\beta-1} D_0[\mathbb{G}(s), \hat{0}]\,ds, \\
& \leq \frac{2 \rho L\mathfrak{D}_{\nu _{2},\nu _{1}}^{\beta}}{\Gamma(\beta+1)}.
\end{align*}
where we use the estimate below via using the assumption (A2)
\begin{equation}\label{eq33}
D_0[\mathbb{G}(\nu ), \hat{0}] \leq 2 L D_{*}[ w, \hat{0}]
=2 \rho L, \quad \forall \nu \in\mathcal{J} 
\end{equation}
Because $\psi$ belongs to $\mathcal{K}$, the distance between 
$(\mathcal{Q} w)(\nu _2)$ and $(\mathcal{Q} w)(\nu _1)$ vanishes as
$\nu _2 \to \nu _1$, which implies the equi-continuity of $\mathcal{Q}$. 
From the Ascoli-Arzelà theorem, it follows that the set 
$\mathcal{Q}(\mathbb{B}(\hat{0}, \rho))$ is relatively compact in 
$\mathcal{X}^{(ii)}$. This analysis shows that $\mathcal{Q}$ is a continuous and 
completely continuous operator.
\smallskip

\noindent\textbf{Step 4.} The final step in invoking Theorem \ref{theo3} consists 
of establishing the boundedness of the set defined by
$$
\mathfrak{P}(\mathcal{Q})=\{ w \in \mathbb{B}(\hat{0}, \rho) :
 u=\lambda \mathcal{Q}(u), \lambda \in(0,1)\},
$$
is bounded in $\mathcal{X}^{(ii)}$. 
Let $ w \in \mathfrak{P}(\mathcal{Q})$ and $\lambda \in(0,1)$, then one notices 
$ w=\lambda \mathcal{Q}(u)$. Hence, we have that
\begin{equation}\label{eq34}
 w(\nu )=\lambda(\varrho(0) \ominus \frac{(-1)}{\Gamma(\beta)} 
 \int_0^{\nu } \psi'(s)\mathfrak{D}_{\nu ,s}^{\beta-1} \mathbb{G}(s)\,ds),
 \quad \nu \in\mathcal{J} 
\end{equation}
By (A3), we obtain
\begin{equation}\label{eq35}
D_0[\mathbb{G}(\nu ), \hat{0}] \leq(\hat{M}_{1}+\hat{M}_{2}) D_{*}[ w, \hat{0}],
 \quad \forall \nu \in\mathcal{J} 
\end{equation}
A straightforward combination of \eqref{eq34} and \eqref{eq35} leads to
$$
D_0[ w(\nu ), \hat{0}] \leq \lambda(D_0[\varrho(0), \hat{0}]+(\hat{M}_{1}
+\hat{M}_{2}) \frac{\mathfrak{D}_{\nu ,0}^{\beta}}{\Gamma(\beta+1)}) D_{*}[ w, \hat{0}],
$$
which implies
$$
D_{*}[ w, \hat{0}] \leq \lambda(D_0[\varrho(0), \hat{0}]+(\hat{M}_{1}+\hat{M}_{2}) 
\frac{\mathfrak{D}_{\nu ,0}^{\beta}}{\Gamma(\beta+1)}) D_{*}[ w, \hat{0}],
$$
where $\psi$ is nondecreasing and $\exists T^{*} \in \mathcal{J}$ such that 
$\lambda(\hat{M}_{1}+\hat{M}_{2})\mathfrak{D}_{T^{*},0}^{\beta}<\Gamma(\beta+1)$, 
where $\beta \in(0,1)$ and $\lambda \in(0,1)$. Then, we obtain
$$
D_{*}[ w, \hat{0}] \leq \lambda D_0[\varrho(0), \hat{0}](1-\lambda(\hat{M}_{1}
+\hat{M}_{2}) \frac{\mathfrak{D}_{T^{*},0}^{\beta}}{\Gamma(\beta+1)})^{-1}<\infty.
$$
The boundedness of $\mathfrak{P}(\mathcal{Q})$ implies, 
via Theorem~\ref{theo3} (Schaefer), that \eqref{sys11} possesses at least one 
solution in $\mathcal{J}^{*}=[0, T^{*}]$.
\end{proof}

\begin{theorem} \label{theo32}
Under assumptions {\rm (A1)--(A3)}, system \eqref{sys11} admits a unique solution 
on $\mathcal{J}^{*}$, $T^{*} \in (0, T]$, for all cases.
\end{theorem}

\begin{proof}
Assume $T^{*} > r$ with the partition
\[
\mathcal{J}^{*} = \cup_{k=0}^{n} [k r, (k+1) r] \cup [(n+1) r, T^{*}],
\]
where $(n+1) r < T^{*} \leq (n+2) r$. For any two solutions $ w$ and 
$\mathfrak{g}$ of \eqref{sys11} satisfying $ w(\nu ) = \varrho_{1}(\nu )$ and 
$\mathfrak{g}(\nu ) = \varrho_{2}(\nu )$ for all $\nu \in [- r, 0]$, 
Lemma~\ref{lem1} yields the following:

For $\nu \in[0, r]$,
\begin{align*}
D_0[ w(\nu ), \mathfrak{g}(\nu )] 
&\leq  D_0[ w(0), \mathfrak{g}(0)]+\frac{ L}{\Gamma(\beta)}
 \int_0^{\nu } \psi'(s)\mathfrak{D}_{\nu ,s}^{\beta-1} D_0[ w(s- r), 
 \mathfrak{g}(s- r)]\,ds \\
&\quad +\frac{ L}{\Gamma(\beta)} \int_0^{\nu } \psi'(s)\mathfrak{D}_{\nu ,s}^{\beta-1} 
D_0[ w(s), \mathfrak{g}(s)]\,ds.
\end{align*}
Considering $ r \in (0, r_{\max}]$, we note that
\[
D_0[ w(\nu ), \mathfrak{g}(\nu )] 
\leq  D_0[ w(0), \mathfrak{g}(0)]+\frac{ L\mathfrak{D}_{\nu ,0}^{\beta}}
{\Gamma(\beta+1)} D_{ r}[\varrho_{1}, \varrho_{2}] 
 +\frac{ L}{\Gamma(\beta)} \int_0^{\nu } \psi'(s)\mathfrak{D}_{\nu ,s}^{\beta-1} 
D_0[ w(s), \mathfrak{g}(s)]\,ds .
\]
Using the generalized Gronwall inequality in \cite{da}, we obtain
\begin{equation}\label{ineq36}
D_0[ w(\nu ), \mathfrak{g}(\nu )] \leq \mathbb{A}_{1}( r) E_{\beta, 1}
( L\mathfrak{D}_{ r,0}^{\beta}), \quad \forall \nu \in[0, r], 
\end{equation}
where
\begin{equation}\label{eq37}
\mathbb{A}_{1}( r)=D_0[ w(0), \mathfrak{g}(0)]
+\frac{ L\mathfrak{D}_{ r,0}^{\beta}}{\Gamma(\beta+1)} D_{ r}[\varrho_{1}, 
\varrho_{2}] . \
\end{equation}

For $\nu \in ( r, 2 r]$, combining \eqref{ineq36} with the inequality 
$\lambda_{1}^{\beta} - \lambda_{2}^{\beta} \leq (\lambda_{1} 
- \lambda_{2})^{\beta}$ (valid for $0 \leq \lambda_{2} \leq \lambda_{1}$, 
$\beta \in (0,1)$) yields
\begin{align*}
D_0[ w(\nu ), \mathfrak{g}(\nu )] 
&\leq  D_0[ w(0), \mathfrak{g}(0)]+\frac{ L\mathfrak{D}_{ r,0}^{\beta}}{\Gamma(\beta+1)}
  D_{ r}[\varrho_{1}, \varrho_{2}] +\mathbb{A}_{1}( r) 
  \frac{ L\mathfrak{D}_{\nu , r}^{\beta}}{\Gamma(\beta+1)} E_{\beta, 1}
  ( L\mathfrak{D}_{ r,0}^{\beta}) \\
&\quad +\frac{ L}{\Gamma(\beta)} \int_0^{\nu } \psi'(s)\mathfrak{D}_{\nu ,s}^{\beta-1}
 D_0[ w(s), \mathfrak{g}(s)]\,ds,
\end{align*}
since $(s- r) \in(- r, 0]$ and $(s- r) \in(0, r]$ for $s \in(0, r]$ and 
$s \in( r, 2 r]$, respectively. Through the generalized Gronwall inequality, we 
establish the following crucial estimate
\begin{equation}\label{ineq38}
D_0[ w(\nu ), \mathfrak{g}(\nu )] \leq \mathbb{A}_{2}( r)
 E_{\beta, 1}( L\mathfrak{D}_{2 r,0}^{\beta}), \quad \forall \nu \in( r, 2 r],
\end{equation}
where
\begin{equation}\label{eq39}
\mathbb{A}_{2}( r)=D_0[ w(0), \mathfrak{g}(0)]
+\frac{ L\mathfrak{D}_{ r,0}^{\beta}}{\Gamma(\beta+1)} D_{ r}[\varrho_{1}, 
\varrho_{2}]+\mathbb{A}_{1}( r) 
\frac{ L\mathfrak{D}_{2 r, r}^{\beta}}{\Gamma(\beta+1)} E_{\beta, 1}
( L\mathfrak{D}_{ r,0}^{\beta}) . 
\end{equation}

By induction on $k \geq 3$, assume $\forall \nu \in ((k-1) r, k r]$ the inequality holds
\begin{equation}\label{ineq310}
D_0[ w(\nu ), \mathfrak{g}(\nu )] \leq \mathbb{A}_{k}( r) 
E_{\beta, 1}( L\mathfrak{D}_{k r,0}^{\beta}), \quad \forall \nu \in((k-1) r, k r], 
\end{equation}
where
\begin{equation}\label{eq311}
\mathbb{A}_{k}( r)=  D_0[ w(0), \mathfrak{g}(0)]
+\frac{ L\mathfrak{D}_{ r,0}^{\beta}}{\Gamma(\beta+1)} D_{ r}[\varrho_{1}, 
\varrho_{2}] 
 +\frac{ L}{\Gamma(\beta+1)} \sum_{i=1}^{k-1}\mathfrak{D}_{(i+1) r,i r}^{\beta} 
\mathbb{A}_{i}( r) E_{\beta, 1}( L\mathfrak{D}_{i r,0}^{\beta}) . 
\end{equation}

For $\nu \in (k r,(k+1) r]$ ($k \geq 3$), using the inequality 
$\mathrm{a}^\beta - \mathrm{b}^\beta \leq (\mathrm{a}-\mathrm{b})^\beta$ with 
$ r \in [0,r_0]$ and \eqref{ineq310}, we obtain
\begin{align*}
&D_0[ w(\nu ), \mathfrak{g}(\nu )] \\
&\leq  D_0[ w(0), \mathfrak{g}(0)]+\frac{ L}{\Gamma(\beta)} 
\sum_{i=0}^{k-1} \int_{i r}^{(i+1) r} \psi'(s)\mathfrak{D}_{\nu ,s}^{\beta-1} 
 D_0[ w(s- r), \mathfrak{g}(s- r)]\,ds \\
&\quad +\frac{ L}{\Gamma(\beta)} \int_{k r}^{\nu } \psi'(s)
 \mathfrak{D}_{\nu ,s}^{\beta-1} D_0[ w(s- r), \mathfrak{g}(s- r)]\,ds 
 +\frac{ L}{\Gamma(\beta)} \int_0^{\nu } \psi'(s)\mathfrak{D}_{\nu ,s}^{\beta-1}
  D_0[ w(s), \mathfrak{g}(s)]\,ds, \\
&\leq  D_0[ w(0), \mathfrak{g}(0)]+\frac{ L\mathfrak{D}_{ r,0}^{\beta}}{\Gamma(\beta+1)} 
 D_{ r}[\varrho_{1}, \varrho_{2}] 
 +\frac{ L}{\Gamma(\beta+1)} \sum_{i=1}^{k-1}\mathfrak{D}_{(i+1) r,i r}^{\beta} 
 \mathbb{A}_{i}( r) E_{\beta, 1}( L\mathfrak{D}_{i r,0}^{\beta}) \\
&\quad +\frac{ L}{\Gamma(\beta+1)}(\psi(\nu )-\psi(k r))^{\beta} \mathbb{A}_{k}
 ( r) E_{\beta, 1}( L\mathfrak{D}_{k r,0}^{\beta})
 +\frac{ L}{\Gamma(\beta)} \int_0^{\nu } \psi'(s)\mathfrak{D}_{\nu ,s}^{\beta-1} 
 D_0[ w(s), \mathfrak{g}(s)]\,ds.
\end{align*}
It follows from the generalized Gronwall inequality that
\begin{equation}\label{ineq312}
D_0[ w(\nu ), \mathfrak{g}(\nu )] 
\leq \mathbb{A}_{k+1}( r) E_{\beta, 1}( L\mathfrak{D}_{(k+1) r,0}^{\beta}), 
\quad \forall \nu \in(k r,(k+1) r] 
\end{equation}
where
\begin{equation} \label{eq313}
\mathbb{A}_{k+1}( r)=  D_0[ w(0), \mathfrak{g}(0)]
+\frac{ L\mathfrak{D}_{ r,0}^{\beta}}{\Gamma(\beta+1)} D_{ r}[\varrho_{1},
\varrho_{2}] +\frac{ L}{\Gamma(\beta+1)}
 \sum_{i=1}^{k}\mathfrak{D}_{(i+1) r,i r}^{\beta} \mathbb{A}_{i}( r) E_{\beta, 1}
 ( L\mathfrak{D}_{i r,0}^{\beta}) .
\end{equation}

 Similarly, for $\nu \in((n+1) r, T^{*}]$, we also have
\begin{align*}
&D_0[ w(\nu ), \mathfrak{g}(\nu )] \\
&\leq  D_0[ w(0), \mathfrak{g}(0)]+\frac{ L}{\Gamma(\beta)} \sum_{i=0}^{n}
  \int_{i r}^{(i+1) r} \psi'(s)\mathfrak{D}_{\nu ,s}^{\beta-1} D_0[ w(s- r), 
  \mathfrak{g}(s- r)]\,ds \\
&\quad +\frac{ L}{\Gamma(\beta)} \int_{(n+1) r}^{\nu } 
 \psi'(s)\mathfrak{D}_{\nu ,s}^{\beta-1} D_0[ w(s- r), \mathfrak{g}(s- r)]\,ds 
 +\frac{ L}{\Gamma(\beta)} \int_0^{\nu } \psi'(s)\mathfrak{D}_{\nu ,s}^{\beta-1} D_0[ w(s), \mathfrak{g}(s)]\,ds, \\
&\leq  D_0[ w(0), \mathfrak{g}(0)]+\frac{ L\mathfrak{D}_{ r,0}^{\beta}}{\Gamma(\beta+1)}
 D_{ r}[\varrho_{1}, \varrho_{2}] 
+\frac{ L}{\Gamma(\beta+1)} \sum_{i=1}^{n}(\psi((i+1) r-\psi(i r))^{\beta} \mathbb{A}_{i}( r) E_{\beta, 1}( L\mathfrak{D}_{i r,0}^{\beta}) \\
&\quad +\frac{ L}{\Gamma(\beta+1)}\mathfrak{D}_{T^{*},(n+1) r}^{\beta} \mathbb{A}_{n+1}
 ( r) E_{\beta, 1}( L\mathfrak{D}_{(n+1) r,0}^{\beta}) 
 +\frac{ L}{\Gamma(\beta)} \int_0^{\nu } \psi'(s)\mathfrak{D}_{\nu ,s}^{\beta-1} D_0[ w(s), \mathfrak{g}(s)]\,ds,
\end{align*}
with $\mathbb{A}_{n+1}( r)$ given by \eqref{eq313}. An application of the generalized 
Gronwall inequality
\begin{equation}\label{ineq314}
D_0[ w(\nu ), \mathfrak{g}(\nu )] \leq \mathbb{A}_{T^{*}}( r)
 E_{\beta, 1}( L\mathfrak{D}_{T^{*},0}^{\beta}), \quad 
 \forall \nu \in((n+1) r, T^{*}], 
\end{equation}
where
\begin{align*}
\mathbb{A}_{T^{*}}( r)
&=  D_0[ w(0), \mathfrak{g}(0)]
 +\frac{ L\mathfrak{D}_{ r,0}^{\beta}}{\Gamma(\beta+1)} 
 D_{ r}[\varrho_{1}, \varrho_{2}] \\
&\quad+\frac{ L}{\Gamma(\beta+1)} \sum_{i=1}^{n}\mathfrak{D}_{(i+1) r,i r}^{\beta} 
 \mathbb{A}_{i}( r) E_{\beta, 1}( L\mathfrak{D}_{i r,0}^{\beta}) \\
&\quad +\frac{ L}{\Gamma(\beta+1)}\mathfrak{D}_{T^{*},(n+1) r}^{\beta} 
 \mathbb{A}_{n+1}( r) E_{\beta, 1}( L\mathfrak{D}_{(n+1) r,0}^{\beta}) .
\end{align*}
Assuming $\varrho_1(\nu ) = \varrho_2(\nu )$ for all $\nu \in [- r, 0]$,
 we observe that
\[
\mathbb{A}_1( r) = \mathbb{A}_2( r) = \dots 
= \mathbb{A}_{n+1}( r) = \mathbb{A}_{T^*}( r) = 0.
\]
Consequently, we obtain the identity $ w(\nu ) = \mathfrak{g}(\nu )$ for all 
$\nu \in [- r, T^*]$.
\end{proof}

\section{Finite time stability result}

\begin{definition} \rm
Problem \eqref{sys11} is said to exhibit FTS with respect to the set 
$\{0,\mathcal{J},\delta,\varepsilon, r\}$ if the following conditions are 
satisfied
\begin{itemize}
 \item[(1)] For any initial function $\varrho$ such that
 $\sup_{\nu \in [- r,0]} D_0[\varrho(\nu ), \hat{0}] \leq \delta$,

\item[(2)] The corresponding solution $ w$ satisfies
 $D_0[ w(\nu ), \hat{0}] \leq \varepsilon \quad \text{for all } \nu \in \mathcal{J}$.
\end{itemize}
\end{definition}

\begin{theorem}\label{theo33}
Under Assumptions{\rm (A1)--(A3)} there exists a solution to \eqref{sys11} 
on \( \mathcal{J}^{*} \). This solution is finite-time stable with respect to 
\( \{0, \mathcal{J}^{*}, \delta, \varepsilon, r\} \) if 
\begin{equation}\label{ineq315}
\delta E_{\beta, 1}(2 L\mathfrak{D}_{T^{*},0}^{\beta}) \leq \varepsilon .
\end{equation}
\end{theorem}

\begin{proof}
Suppose that \( w \) solves system \eqref{sys11} over the interval
\( \mathcal{J}^{*} \). We define the function \( \mathcal{Z}(\nu ) \) to measure 
the distance between \( w(\nu ) \) and the zero element. 
Then, by applying the explicit solution formula from \eqref{sys28} and using 
assumption (A2), we can establish the following estimate for 
all \( \nu \in \mathcal{J}^{*} \),
\begin{equation}\label{ineq316}
\mathcal{Z}(\nu ) \leq \mathcal{Z}(0)+\frac{1}{\Gamma(\beta)} 
\int_0^{\nu } \psi'(s)\mathfrak{D}_{\nu ,s}^{\beta-1}[ L \mathcal{Z}(s)
+ L \mathcal{Z}(s- r)]\,ds. 
\end{equation}
Let $\mathcal{Z}_{*}(\nu )=\sup _{\xi \in[- r, t]} \mathcal{Z}(\xi)$ , 
for all $\nu \in\mathcal{J}^{*}$. Clearly, \( \mathcal{Z}_{*}(\nu ) \) increases with 
\( t \), and we note that
\begin{equation}\label{ineq317}
\mathcal{Z}(\nu - r) \leq \mathcal{Z}_{*}(\nu ) \quad\text{and}\quad
\mathcal{Z}(\nu ) \leq \mathcal{Z}_{*}(\nu ). 
\end{equation}
Moreover, we observe that \( \mathcal{Z}(s - r) \leq \mathcal{Z}_{*}(s) \) 
for all \( s \in \mathcal{J} \). Consequently, inequality \eqref{ineq316} implies 
that
\begin{equation}\label{ineq318}
\mathcal{Z}(\nu ) \leq \mathcal{Z}(0)+\frac{1}{\Gamma(\beta)} 
\int_0^{\nu } \psi'(s)\mathfrak{D}_{\nu ,s}^{\beta-1}2 L \mathcal{Z}_{*}(s)\,ds .
\end{equation}

Next, we set $C=2 L$ and $\mathcal{U}(\nu )=\frac{C}{\Gamma(\beta)}
 \int_0^{\nu } \psi'(s)\mathfrak{D}_{\nu ,s}^{\beta-1} \mathcal{Z}_{*}(s)\,ds$.
We want to show that \( \mathcal{U}(\nu ) \) gets larger as time \( t \) 
increases within \( \mathcal{J}^{*} \). 
Take two times \( \nu _1 \) and \( \nu _2 \) such that 
\( \nu _1 \leq \nu _2 \). Since
$$
\frac{\mathfrak{D}_{\nu _2,0}^{\beta}}{\Gamma(\beta+1)} 
\geq \frac{\mathfrak{D}_{\nu _1,0}^{\beta}}{\Gamma(\beta+1)},
$$
it follows that
$$
\frac{C}{\Gamma(\beta)} \int_{\nu _{1}}^{\nu _{2}} \psi'(s)
\mathfrak{D}_{\nu _{2},0}^{\beta-1}\,ds 
\geq \frac{C}{\Gamma(\beta)} \int_0^{\nu _{1}}[\psi'(s)\mathfrak{D}_{\nu _{1},0}^{\beta-1} -\psi'(s)\mathfrak{D}_{\nu _{2},0}^{\beta-1} ]\,ds.
$$
Since $\mathcal{Z}_{*}$ is increasing on $\mathcal{J}^{*}$, one has
$$
\frac{C}{\Gamma(\beta)} \int_{\nu _{1}}^{\nu _{2}} 
\psi'(s)\mathfrak{D}_{\nu _{2},0}^{\beta-1} \mathcal{Z}_{*}(s)\,ds 
\geq \frac{C}{\Gamma(\beta)} \int_0^{\nu _{1}}[\psi'(s)
\mathfrak{D}_{\nu _{1},0}^{\beta-1} -\psi'(s)\mathfrak{D}_{\nu _{2},0}^{\beta-1} ]
 \mathcal{Z}_{*}(s)\,ds.
$$
Therefore, 
\begin{align*}
&\mathcal{U}(\nu _{2})-\mathcal{U}(\nu _{1}) \\
& =\frac{C}{\Gamma(\beta)} \int_0^{\nu _{2}} \psi'(s)
 \mathfrak{D}_{\nu _{2},0}^{\beta-1} \mathcal{Z}_{*}(s)\,ds
  -\frac{C}{\Gamma(\beta)} \int_0^{\nu _{1}} \psi'(s)\mathfrak{D}_{\nu _{1},0}^{\beta-1}
   \mathcal{Z}_{*}(s)\,ds, \\
& =\frac{C}{\Gamma(\beta)} \int_{\nu _{1}}^{\nu _{2}} \psi'(s)
 \mathfrak{D}_{\nu _{2},0}^{\beta-1} \mathcal{Z}_{*}(s)\,ds
 -\frac{C}{\Gamma(\beta)} \int_0^{\nu _{1}}[\psi'(s)\mathfrak{D}_{\nu _{1},0}^{\beta-1}
  -\psi'(s)\mathfrak{D}_{\nu _{2},0}^{\beta-1} ] 
 \mathcal{Z}_{*}(s)\,ds \geq 0.
\end{align*}
The above inequality shows that
\[
\mathcal{U}(\nu_2)-\mathcal{U}(\nu_1)\ge 0
\quad \text{for all } \nu_2\ge\nu_1.
\]
This monotonicity follows from  \(\mathcal{Z}_{*}(\nu)\) being increasing,
and the positivity of the generalized fractional kernel
$(\psi(\nu)-\psi(s))^{\beta-1}$, since $\psi\in\mathcal{K}$ is strictly
increasing and $\psi'(s)>0$.
Hence, the generalized fractional integral operator preserves
monotonicity, and $\mathcal{U}(\nu)$ increasing on $\mathcal{J}^*$.
Applying inequality \eqref{ineq318}, we obtain
\begin{equation}\label{ineq320}
\mathcal{Z}(\xi) \leq \mathcal{Z}(0)+\frac{C}{\Gamma(\beta)} 
\int_0^{\nu } \psi'(s)\mathfrak{D}_{\nu ,s}^{\beta-1} \mathcal{Z}_{*}(s)\,ds.
\end{equation}
for all $\xi \in\mathcal{J}$. Thus, we have
\begin{equation}\label{eq321}
\mathcal{Z}_{*}(\nu )=\sup _{\xi \in[- r, t]} \mathcal{Z}(\xi) 
\leq \mathcal{Z}(0)+\frac{C}{\Gamma(\beta)} \int_0^{\nu } 
\psi'(s)\mathfrak{D}_{\nu ,s}^{\beta-1} \mathcal{Z}_{*}(s)\,ds. 
\end{equation}
By applying the generalized Gronwall inequality, we obtain that
\begin{equation}\label{ineq322}
\mathcal{Z}_{*}(\nu ) \leq \mathcal{Z}(0) E_{\beta, 1}(C\mathfrak{D}_{\nu ,0}^{\beta}).
\end{equation}
So, this means that
$$
\mathcal{Z}(\nu ) \leq \mathcal{Z}_{*}(\nu )=\sup _{\xi \in[- r, t]} 
\mathcal{Z}(\xi) \leq \mathcal{Z}(0) E_{\beta, 1}(C\mathfrak{D}_{T^{*},0}^{\beta}),
$$
where $\mathcal{Z}(0)=D_0[\varrho(0), \hat{0}]$. Then, if 
$D_0[\varrho(0), \hat{0}] \leq \delta$, it follows from \eqref{ineq315} that
$$
D_0[ w(\nu ), \hat{0}] \leq \varepsilon, \quad \forall \nu \in\mathcal{J}^{*} .
$$
This yields the FTS of the solution of problem \eqref{sys11}.
\end{proof}

\begin{theorem}\label{theo34}
Under assumptions {\rm (A1)--(A3)}there exists a solution to \eqref{sys11} on 
\( \mathcal{J}^{*} \). This solution is finite-time stable with respect to 
\( \{0, \mathcal{J}^{*}, \delta, \varepsilon, r\} \) if 
\begin{equation}\label{ineq323}
\mathbb{E}_{T^{*}}( r) E_{\beta, 1}( L\mathfrak{D}_{T^{*},0}^{\beta}) 
\leq \varepsilon, 
\end{equation}
where $\mathbb{E}_{T^{*}}( r)$ is defined in the proof below.
\end{theorem}

\begin{proof}
Assume that \( T > r \). We partition the interval \( \mathcal{J} \) as
\[
\mathcal{J} = \cup_{k=0}^{n} [k r, (k+1) r] \cup [(n+1) r, T],
\]
where \( (n+1) r < T \leq (n+2) r \). By applying the method of steps in conjunction 
with the generalized Gronwall inequality, we establish the FTS of system \eqref{sys11}, 
as outlined below. Additionally, we assume that the initial condition satisfies 
\( D_0[\varrho(\nu ), \hat{0}] \leq \delta \) for all \( \nu \in [- r, 0] \).

Consider $\nu \in [0, r]$. Applying the same estimation method developed for inequality 
\eqref{ineq36} in the proof of Theorem~\ref{theo32}, and noting that $ r \leq r_{\max}$, 
we establish
\begin{equation}\label{ineq324}
D_0[ w(\nu ), \hat{0}] \leq \mathbb{E}_{1}( r) E_{\beta, 1}
( L\mathfrak{D}_{ r,0}^{\beta}), \quad \forall \nu \in[0, r], \tag{3.24}
\end{equation}
where
\begin{equation}\label{eq325}
\mathbb{E}_{1}( r)=\delta(1+\frac{ L\mathfrak{D}_{ r,0}^{\beta}}{\Gamma(\beta+1)}) .
\end{equation}

 For $\nu \in( r, 2 r]$, as in the estimate \eqref{ineq38}, we obtain
\begin{equation}\label{ineq326}
D_0[ w(\nu ), \hat{0}] \leq \mathbb{E}_{2}( r) E_{\beta, 1}
( L\mathfrak{D}_{2 r,0}^{\beta}), \quad \forall \nu \in( r, 2 r], 
\end{equation}
where
$$
\mathbb{E}_{2}( r)=\delta(1+\frac{ L\mathfrak{D}_{ r,0}^{\beta}}{\Gamma(\beta+1)})
+\mathbb{E}_{1}( r) \frac{ L\mathfrak{D}_{2 r, r}^{\beta}}{\Gamma(\beta+1)} E_{\beta, 1}
( L\mathfrak{D}_{ r,0}^{\beta}).
$$

For $\nu \in(k r,(k+1) r]$, where $k=3,4,5, \ldots$. Similar to the estimate 
\eqref{ineq312}, the following estimate holds,
\begin{equation}\label{ineq327}
D_0[ w(\nu ), \hat{0}] \leq \mathbb{E}_{k+1}( r) E_{\beta, 1}
( L\mathfrak{D}_{(k+1) r,0}^{\beta}), \quad \forall \nu \in(k r,(k+1) r] 
\end{equation}
where
$$
\mathbb{E}_{k+1}( r)
=  \delta(1+\frac{ L\mathfrak{D}_{ r,0}^{\beta}}{\Gamma(\beta+1)})
 +\frac{ L}{\Gamma(\beta+1)} \sum_{i=1}^{k}\mathfrak{D}_{(i+1) r,i r}^{\beta}
  \mathbb{E}_{i}( r) E_{\beta, 1}( L\mathfrak{D}_{i r,0}^{\beta}).
$$

 For $\nu \in ((k+1) r, T^{*}]$, following the same bounding technique as in 
 \eqref{ineq314}, we establish
\[
D_0[ w(\nu ), \hat{0}] \leq \varepsilon,
\]
which verifies the FTS of solution to problem \eqref{sys47}.
\begin{equation}\label{ineq328}
D_0[ w(\nu ), \hat{0}] 
\leq \mathbb{E}_{T^{*}}( r) E_{\beta, 1}( L\mathfrak{D}_{T^{*},0}^{\beta}), \quad 
\forall \nu \in((n+1) r, T^{*}], 
\end{equation}
where
\begin{align*}
\mathbb{E}_{T^{*}}( r)
&=  \delta(1+\frac{ L\mathfrak{D}_{ r,0}^{\beta}}{\Gamma(\beta+1)})
+ \frac{ L}{\Gamma(\beta+1)} \sum_{i=1}^{n}\mathfrak{D}_{(i+1) r,i r}^{\beta}
 \mathbb{E}_{i}( r) E_{\beta, 1}( L\mathfrak{D}_{i r,0}^{\beta}) \\
&\quad +  \frac{ L}{\Gamma(\beta+1)}\mathfrak{D}_{T^{*},(n+1) r}^{\beta} 
\mathbb{E}_{n+1}( r) E_{\beta, 1}( L\mathfrak{D}_{(n+1) r,0}^{\beta}) .
\end{align*}


For $\nu \in ((k+1) r, T^{*}]$, akin to the bound in \eqref{ineq314}, we derive 
$D_0[ w(\nu ), \hat{0}] \leq \varepsilon$, which ensures the FTS of the solution 
to system \eqref{sys47}.
\end{proof}

\begin{remark}\label{rem31} \rm
Using the analytical technique from Theorem \ref{theo34} applied to the interval 
\( T^{*} \in (0, r] \), we arrive at the  estimate
$$
D_0[ w(\nu ), \hat{0}] 
\leq \delta(1+\frac{ L\mathfrak{D}_{ r,0}^{\beta}}{\Gamma(\beta+1)}) E_{\beta, 1}
( L\mathfrak{D}_{ r,0}^{\beta}) .
$$
Consequently, system \eqref{sys11} exhibits FTS under the following sufficient 
condition.
\begin{equation}\label{ineq329}
\delta(1+\frac{ L\mathfrak{D}_{ r,0}^{\beta}}{\Gamma(\beta+1)}) E_{\beta, 1}
( L\mathfrak{D}_{ r,0}^{\beta}) \leq \varepsilon .
\end{equation}
\end{remark}

\section{Examples}

In this section, we illustrate the validate of our theoretical results. 
While the computational framework developed in Theorem~\ref{theo34} appears more 
involved than that of Theorem~\ref{theo33}, we emphasize that its FTS conditions 
are actually less restrictive, offering broader applicability in practice.
We examine the concrete problem
\begin{equation} \label{sys47}
\begin{gathered}
{ }^C \mathcal{D}_{0^{+}}^{\beta ; \psi} w(\nu )=\mathcal{H}(\nu , w(\nu ),
w(\nu - r)), \quad \nu \in[0,0.3], \\
 w(\nu )=\varrho(\nu ) \in \mathfrak{F}, \quad \nu \in[- r, 0],
\end{gathered}
\end{equation}
For the function $\mathcal{H}(\nu , w(\nu ), w(\nu - r)) = \sin(\nu ) w(\nu )
+ \cos(\nu ) w(\nu - r)$, where the delay is constant at $ r = 0.1$, we suppose
$\sup_{\nu \in [- r, 0]} D_0[\varrho(\nu ), \hat{0}] \leq \delta$, with $\delta$
being a positive real number. Due to the difficulty in deriving an analytical
solution for \eqref{sys47}, this example provides sufficient conditions for the FTS
of the solution, setting $L = 1$. The solution is presumed to exist over the interval
$[0, 0.3]$.

 From Theorem \ref{theo33}, the solution to \eqref{sys47} is FTS with respect to 
 $ \{0, [0, 0.3], \delta, \varepsilon, 0.1\}$ provided the specified conditions hold.
$$
E_{\beta, 1}(2\mathfrak{D}_{0.3,0}^{\beta}) \leq \varepsilon / \delta .
$$

According to Theorem \ref{theo34}, the solution of system \eqref{sys47} demonstrates 
FTS with respect to $\{0, [0, 0.3], \delta, \varepsilon, 0.1\}$ when the given 
conditions are satisfied.
$$
\mathbb{E}_{0.3}(0.1) E_{\beta, 1}(\mathfrak{D}_{0.3,0}^{\beta}) \leq \varepsilon,
$$
where
\begin{gather*}
 \mathbb{E}_{1}(0.1)=\delta(1+\frac{\mathfrak{D}_{0.1,0}^{\beta}}{\Gamma(\beta+1)}), \\
 \mathbb{E}_{2}(0.1)=\mathbb{E}_{1}(0.1)+\mathbb{E}_{1}(0.1) \frac{\mathfrak{D}_{0.2,0.1}^{\beta}}{\Gamma(\beta+1)} E_{\beta, 1}(\mathfrak{D}_{0.1,0}^{\beta}), \\
 \mathbb{E}_{0.3}(0.1)=\mathbb{E}_{2}(0.1)+\mathbb{E}_{2}(0.1) \frac{\mathfrak{D}_{0.3,0.2}^{\beta}}{\Gamma(\beta+1)} E_{\beta, 1}(\mathfrak{D}_{0.2,0}^{\beta}).
\end{gather*}


\begin{table}[htb]
\arraystretch{1.2}
\centering
\caption{Evaluation of Theorems \ref{theo33} and \ref{theo34} for fuzzy time-scale 
stability with $\delta = 20$.}
\label{table1}
\begin{tabular}{|lccc|}
\hline
\textbf{Sufficient conditions on $\psi$} & $t$ & $\log (t+1)$ & $\sqrt{t^{2}+1}$ \\
\hline
Theorem \ref{theo33} & $1.247477 \times 10^{2}$ & $1.058240 \times 10^{2}$ & $34.51709$ \\
Theorem \ref{theo34} & $2.6871 \times 10^{3}$ & $2.1038 \times 10^{3}$ & $161.9795$ \\
\hline
\end{tabular}
\end{table}

To illustrate the applicability of Theorems \ref{theo33} and \ref{theo34} 
in analyzing the FTS of system \eqref{sys47}, Table \ref{table1} presents
numerical results  for different functions $\psi$, with $\beta = 0.5$.

\section*{Conclusion}
 In this article, we studied fuzzy fractional delay differential equations within 
the framework of the generalized Caputo fractional derivative. By employing stepwise 
approximation techniques together with Gronwall-type inequalities, we established 
the existence and uniqueness of solutions. Furthermore, sufficient conditions 
ensuring finite-time stability were derived. Theoretical findings were supported 
with computational simulations, which validated the analytical results and 
demonstrated the effectiveness of the proposed framework. 
The obtained results highlight the interplay between fuzziness, memory effects, 
and delays in fractional dynamical systems, providing new insights into their 
qualitative behavior.

\subsection*{Acknowledgements}
We are very grateful to the anonymous referee, who helped us improve the article.

\begin{thebibliography}{10}

\bibitem{abbas}  Abbas, A.; Mubeen, S.; Sohai, T.; Hussain, S.; Budak, H.;
\emph{Exploring Milne-type inequalities through the generalized Riemann-Liouville 
fractional operator}, Filomat, 39 (24), 8575-8605, (2025).
 
\bibitem{ahm}
Ahmad, S.; Ullah, A.; Hoa, N. V.;
\emph{Fuzzy natural transform method for solving fuzzy differential equation}.
Soft Computing, 27 (2013), 8611-8625.

\bibitem{ahm1} Ahmed, H. M.; Ragusa, M. A.;
\emph{Nonlocal controllability of Sobolev-type conformable fractional stochastic 
evolution inclusions with Clarke subdifferential}, 
Bulletin of the Malaysian Mathematical Sciences Society, 45 (2022) (6), 3239-3253, 
doi: 10.1007/s40840-022-01377-y.

\bibitem{akd} Akdemir, A. O.; Aslan, S.; Dokuyucu, M.; Celik, E.;
\emph{Exponentially convex functions on the coordinates and novel estimations via
 Riemann-Liouville fractional operator}, Journal of Function Spaces, vol. 2023, (2023).


\bibitem{allah2014}
Allahviranloo, T.; Gouyandeh, Z.; Armand, A.;
\emph{Fuzzy fractional differential equations under generalized fuzzy Caputo derivative}.
Journal of Intelligent \& Fuzzy Systems, 26  (2014), 1481-1490.

\bibitem{alm}
Almeida, R.;
\emph{A Caputo fractional derivative of a function with respect to another
function}, Communications in Nonlinear Science and Numerical Simulation, 44 (2017),
 460-481.
 
\bibitem{bal}
Baleanu, D.; Shiri, B.;
\emph{Generalized fractional differential equations for past dynamic}.
AIMS Mathematics, 7 (2022), 14394-14418.

\bibitem{bed}
Bede, B.; Stefanini, L.;
\emph{Generalized differentiability of fuzzy-valued functions}. 
\emph{Fuzzy Sets and Systems}, 230 (2013), 119-141.

\bibitem{che}
Chen, L.; Wu, W. P. R.;  He, Y.;
\emph{New result on FTS of fractional-order nonlinear delayed systems}. 
Journal of Computational and Nonlinear Dynamics, 10 (2015) (6), 064504.

\bibitem{ali2} El Mfadel, A.; Melliani, S.; Elomari, M. H.;
\emph{Notes on local and nonlocal intuitionistic fuzzy fractional boundary
value problems with Caputo fractional derivatives}.
Journal of Mathematics, 2021 (2021) (1), 4322841.

\bibitem{ali3} El Mfadel, A.; Melliani, S.; Elomari, M. H.;
\emph{On the existence and uniqueness results for fuzzy linear and semilinear
fractional evolution equations involving Caputo fractional derivative}.
Journal of Function Spaces, 2021 (2021) (1), 4099173.

\bibitem{Fan}
Fan, Q.; Huang, L. L.;Wu, G. C.;
\emph{ General fractional interval-valued differential equations and Gronwall
inequalities}. \emph{Soft Computing}, 27 (2013), 7739-7749.

\bibitem{hoa2019b}
Hoa, N. V.; Vu, H.;
\emph{A survey on the initial value problems of fuzzy implicit fractional
differential equations}. Fuzzy Sets and Systems,  400 (2020), 90-133.

\bibitem{hoa2020a}
Hoa, N. V.;
\emph{On the initial value problem for fuzzy differential equations of
non-integer order $\beta \in(1,2)$}. Soft Computing, 24 (2020), 935-954.

\bibitem{hoa2018}
Hoa, N. V.; Lupulescu, V.; O'Regan, D.;
\emph{A note on initial value problems for fractional fuzzy
differential equations}. Fuzzy Sets and Systems, 347 (2018), 54-69.


\bibitem{huan2022}
Huang, L. L.; Wu, G. C.; Fan, Q.; Shiri, B.;
\emph{Fractional linear interval-valued systems with w-monotonicity's
constraint conditions}. Fuzzy Systems and Mathematics. (2022)

\bibitem{laz2006}
Lazarevi\'c, M. P.;
\emph{Finite time stability analysis of PD fractional control of robotic time-delay
systems}. Mechanics Research Communications, 33 (2006), 269-279.

\bibitem{long}
Long, H. V.; Son, N. T. K.; Tam, H. T. T.;
\emph{The solvability of fuzzy fractional partial differential equations
under Caputo gH-differentiability}. Fuzzy Sets and Systems, 309 (2017), 35-63.

\bibitem{luo}
Luo, C.; Wu, G. C.;Huang, L. L.;
\emph{Fractional uncertain differential equations with general memory effects:
Existences and alpha-path solutions}.
Nonlinear Analysis: Modelling and Control, 28 (2023) (1), 1-28.

\bibitem{lupu}
Lupulescu, V.;
\emph{Fractional calculus for interval-valued functions}.
Fuzzy Sets and Systems, 265 (2015), 63-85.

\bibitem{ali1} 
Melliani, S.; El Mfadel, A. E.; Chadli, L. S.; Elomari, M. H.;
\emph{Stability of intuitionistic fuzzy nonlinear fractional differential 
equations}. Notes Instuitionistic Fuzzy Sets, 27 (2021) (1), 83-100.

\bibitem{naif}
Naifar, O.; Nagy, A. M.; Makhlouf, A. B.; Kharrat, M.;  Hammami, M. A.;
\emph{FTS of linear fractional-order time-delay systems}.
International Journal of Robust and Nonlinear Control, 29 (2019), 180-187.

\bibitem{phat}
Phat, V. N.; Thanh, N. T.;
\emph{New criteria for finite-time stability of nonlinear fractional-order
delay systems: A Gronwall inequality approach}.
Applied Mathematics Letters, 83 (2018), 169-175.

\bibitem{sal}
Salehi, O. S.; Salehi, M.;
\emph{A survey on fuzzy fractional variational problems}. 
Journal of Computational and Applied Mathematics, 271 (2014), 71-82.

\bibitem{samk}
Samko, S. G.; Kilbas, A. A.; Marichev, O. I.;
\emph{Fractional integrals and derivatives: Theory and applications}.
Gordon and Breach, 1993.

\bibitem{smar}
Smart, D. R.;
\emph{Fixed point theorems} (Vol. 66). Cambridge University Press, 1980

\bibitem{da}
da Sousa, J. V.; de Oliveira, C.; de Oliveira, E. C.;
\emph{A Gronwall inequality and the Cauchy-type problem by means of 
$\psi$-Hilfer operator}. 
\emph{Differential Equations and Applications}, 11 (2019), 87-106.

\bibitem{shiri}
Wu, G. C.; Shiri, B.; Fan, Q.; Feng, H. R.;
\emph{Terminal value problems of non-homogeneous fractional linear systems with 
general memory kernels}. Journal of Nonlinear Mathematical Physics, 30 (202), 303-314.

\bibitem{xi}
Xiang, C.; Duan, J.; Luo, Q.;
\emph{Finite time blow-up for an inhomogeneous wave equation with Riemann-Liouville
fractional integral}, Electronic Journal of Applied Mathematics, 3 (2), (2025).

\bibitem{zhan}
Zhang, F.; Qian, D.; Li, C.;
\emph{Finite-time stability analysis of fractional differential systems 
with variable coefficients}. Chaos, 29 (2019), 013110.

\end{thebibliography}

\end{document}

